\label{chap:p3-83-27-menos_un_doceavo}

\section{Hensel--Witt domain, Teichmüller maps, and rational target \(-1/12\)}

The NTHW is the data and maps package
\[
\NTHW=
\bigl(
\mathscr R_\pi^{\rm micro},
\Gnine,
\Cdoce,
K,
\mathcal W_{12\to24},
\mathcal I_{6\to8}
\bigr),
\]
with the following compatibilities:
\begin{enumerate}
  \item the Hensel regions share a quotient word and retain
  distinct memory;
  \item \(\Gnine\) transports orientation and produces monodromy;
  \item \(\Cdoce\) organizes the signed observable and the seal \(K\);
  \item the Witt reading converts supports of the same boundary into hexads and
  stars;
  \item the extension \(W_{12}\to W_{24}\) produces the
  hexad--octad interface;
  \item the Archimedean, exceptional and dimensional readers receive
  invariants of the same state.
\end{enumerate}

\begin{figure}[!htbp]
\centering
\resizebox{\textwidth}{!}{\begin{tikzpicture}[font=\small,node distance=1.2cm and 1cm,
  nodebox/.style={draw,rounded corners=2mm,align=center,minimum width=3cm,minimum height=1.05cm},
  arr/.style={-{Stealth[length=2.4mm]},thick}]
  \node[nodebox,draw=hmtblue,fill=hmtlightblue] (hensel) at (-5,2.3) {Henselian cylinders\\carry cocycle and \(\MonNine\)};
  \node[nodebox,draw=hmtblue,fill=hmtlightblue] (regions) at (-5,-.1) {five regions of \(\pi\)\\regions of \(e,\varphi\)};
  \node[nodebox,draw=hmtgold,fill=hmtlightgold,minimum width=3.8cm,minimum height=2.05cm,font=\bfseries] (core) at (0,1.1) {\(\mathscr T_{\rm HW}\)\\Correspondences\\Hensel--Witt};
  \node[nodebox,draw=hmtviolet,fill=hmtlightviolet] (c12) at (0,-1.8) {\(\Cdoce\), \(\Usgn\), \(K\)\\differences \(D_3,D_4,Q\)};
  \node[nodebox,draw=hmtgreen,fill=hmtlightgreen] (witt) at (5,2.3) {Witt \(W_{12}\)\\supports and \(M_{12}\)};
  \node[nodebox,draw=hmtgreen,fill=hmtlightgreen] (oct) at (5,-.1) {incidence inclusion \(6\to8\)\\channels \(90/120\), normalization \(-1/12\)};
  \node[nodebox,draw=hmtred,fill=hmtlightred] (readers) at (5,-2.5) {coordinated readers\\value and symmetry};
  \draw[arr,hmtblue] (hensel)--(core);
  \draw[arr,hmtblue] (regions)--(core);
  \draw[arr,hmtgold] (core)--(c12);
  \draw[arr,hmtgreen] (core)--(witt);
  \draw[arr,hmtgreen] (core)--(oct);
  \draw[arr,hmtred] (c12)--(readers);
  \draw[arr,hmtred] (witt.east)--++(1.9,0)|-(readers.east);
  \draw[arr,hmtred] (oct)--(readers);
  \node[align=center,hmtgray,font=\footnotesize] at (0,-3.45) {Compatibility of cylinders, dodecaphase, and incidences on the same enriched state.};
\end{tikzpicture}
}
\caption{The NTHW as a typed interface. The central mathematical task is to make
the diagram commute, not to accumulate coincidences around a name.}
\label{p3-83-c27-s1:fig:nthw}
\end{figure}

NTHW gives a formal name to the place previously described as the “Rosetta
Stone.” The substitution is conceptual: a stone translates through resemblance; a
transduction specifies input, output, preserved datum and loss of
information.

\section{Teichmüller map from the ternary code to the Witt ring}

The reversible APP boundary selects, with declared gauges, a matrix
\(A_W\) over \(\FF_3\) satisfying
\[
A_WA_W^T=2I\pmod3.
\]
The graphic code
\[
C_W=\{(q,qA_W):q\in\FF_3^6\}
\]
is the extended ternary Golay code. Its weight enumerator is
\[
1+264y^6+440y^9+24y^{12}.
\]
The \(264\) words of weight six, identified up to sign, yield \(132\)
supports. Every five-point subset belongs to a unique hexad:
\[
W_{12}=S(5,6,12).
\]

The field of \(495\) star involutions generates a group of order
\(95\,040\), recognized as \(M_{12}\). An isolated star does not generate the
group. The ordered HMT frame has trivial stabilizer and therefore fixes a
gauge within the action.

\section{Hexad--octad incidence and modular channels \(90/120\)}

Let \(H\) be a hexad and \(O_H\) an octad containing it. Marking one point
and a pair of points of the hexad yields
\[
F_6(H)=H\times\binom H2,
\qquad
|F_6|=6\binom62=90.
\]
If the marked point can traverse the octad:
\[
F_8(H)=O_H\times\binom H2,
\qquad
|F_8|=8\binom62=120.
\]
The inclusion \(H\hookrightarrow O_H\) induces \(F_6\hookrightarrow F_8\).
Its normalized defect is
\[
\boxed{
\frac{90-120}{24\binom62}
=\frac{6-8}{24}
=-\frac1{12}.}
\]
This combinatorial theorem precedes Barbero’s angular reading. The
two objects share the interface \(6\to8\), but are not identical.

\section{4 compatible routes to \texorpdfstring{\(-1/12\)}{-1/12}}

The construction contains four realizations that must be presented together, not
counted as four independent proofs of a single provenance.

\subsection{Angular realization through the elementary increment dodecaphasic}

One tooth of \(\Cdoce\) measures \(30^\circ\). Then
\[
\frac{30}{120}-\frac{30}{90}
=\frac14-\frac13
=-\frac1{12}.
\]
It is the angular shadow of the pattern \(90/120\).

\subsection{Oriented pseudoinverse of the Gram operator on the surviving subspace}

In the certified jump \(60\to81\), two fibers have sizes \(0\) and
\(12\). The Gram operator is
\[
K_{60,81}=\operatorname{diag}(0,12).
\]
On the live subspace, its pseudoinverse is \(1/12\); with the negative
boundary orientation,
\[
E_{\rm surv}=-K_{60,81}^{+}
\]
yields \(-1/12\).

\subsection{Scale extraction by the 2nd triadic difference}

The realization via a triadic scale extractor is the second-difference construction
developed in full in
\cref{sec:c27-segunda-diferencia-triadica}. This first occurrence establishes its
role within the inventory of realizations; the indicated residence preserves
the definitions of \(T_L\), \(a(L)\), \(A_1(L)\), \(C(L)\), the independence
of the residue from \(L\), and the exact status of the normalization.

\subsection{Modular realization through the quotient \(\eta(\tau)/\eta(3\tau)\) and the function \(t_3\)}

The eta function satisfies
\[
\frac{\eta(\tau)}{\eta(3\tau)}
=q^{-1/12}\prod_{n\ge1}(1+q^n+q^{2n})^{-1}.
\]
The exponent comes from \(1/24-3/24=-1/12\). Taking twelve copies:
\[
t_3(q)=\left(\frac{\eta(\tau)}{\eta(3\tau)}\right)^{12}
=q^{-1}-12+54q-76q^2-243q^3+\cdots.
\]
The coefficient \(54\) coincides with the APP defect. The corridor
\[
j=\frac{(t_3+27)(t_3+243)^3}{t_3^3}
\]
connects afterwards with the classical modular branch.

\begin{remark}[What Has Been Proved in the Corridor]
The APP defect \(54\), the eta exponent \(-1/12\), the
hexad--octad defect, and the modular identities once \(t_3\) has been fixed are exact. The
strong assertion is that all of them are natural images of a single TPK operator.
That naturality must be proved by means of the NTHW diagram; it does not
follow merely from repeating the same integers.
\end{remark}


\section{5 algebraic realizations of \texorpdfstring{\(-1/12\)}{-1/12} and separation of their domains}

In addition to the normalized defect of the hexad--octad inclusion,
five operator realizations of the same rational number appear. Their domains are
specified separately; equality of values does not identify the
operators without an additional intertwiner.

\subsection{Difference of dodecaphasic periodicities}

One tooth of \(\Cdoce\) measures \(30^\circ\). Then
\[
\frac{30}{120}-\frac{30}{90}
=\frac14-\frac13
=-\frac1{12}.
\]
It is the oriented difference between the periods of \(90^\circ\) and
\(120^\circ\) in this realization.

\subsection{Pseudoinverse of the Gram Operator on the Live Component}

In the exact jump \(60\to81\), two fibers have sizes \(0\) and
\(12\). The Gram operator is
\[
K_{60,81}=\operatorname{diag}(0,12).
\]
On the nonzero component of its range, the pseudoinverse is \(1/12\); with the negative
boundary orientation,
\[
E_{\rm surv}=-K_{60,81}^{+}
\]
yields \(-1/12\).

\subsection{2nd difference between triadic scales}
\label{sec:c27-segunda-diferencia-triadica}

Define
\[
T_L=\sum_{n=1}^{L-1}n(L-n)=\frac{L(L^2-1)}6,
\qquad
a(L)=-\frac{T_L}{L^2}=-\frac L6+\frac1{6L}.
\]
The first scale subtraction is
\[
A_1(L)=a(L)-\frac{a(3L)}3=\frac4{27L},
\]
and the second
\[
C(L)=LA_1(L)-\frac{3L}{3}A_1(3L)=\frac8{81}.
\]
The residue \(8/81\) is independent of \(L\). To obtain
\(-1/12\), the declared normalization is applied
\[
R(L)=-\frac{27}{32}C(L)=-\frac1{12}.
\]
The calculation derives \(8/81\); it does not by itself derive the factor \(-27/32\).

\subsection{Cadena modular}

The eta function satisfies
\[
\frac{\eta(\tau)}{\eta(3\tau)}
=q^{-1/12}\prod_{n\ge1}(1+q^n+q^{2n})^{-1}.
\]
The exponent comes from \(1/24-3/24=-1/12\). Taking twelve copies:
\[
t_3(q)=\left(\frac{\eta(\tau)}{\eta(3\tau)}\right)^{12}
=q^{-1}-12+54q-76q^2-243q^3+\cdots.
\]
The coefficient \(54\) coincides with the APP defect\@. The rational identity
\[
j=\frac{(t_3+27)(t_3+243)^3}{t_3^3}
\]
establishes the connection with the classical modular branch.

\Needspace{32\baselineskip}
\subsection{Defect of cyclic difference projectors}

Let \(S\) be the cyclic shift on \(\QQ^{12}\) and define
\[
 D_r=I-S^r.
\]
The kernel of \(D_r\) consists of the vectors constant on each orbit
of \(S^r\); therefore,
\[
 \dim\ker D_r=\gcd(12,r).
\]
Let \(\Pi_{\ker D_3}\) and \(\Pi_{\ker D_4}\) be the rational projectors onto
\(\ker D_3\) and \(\ker D_4\), and let
\(\tau(T)=\operatorname{Tr}(T)/12\) be the normalized trace. Then
\[
 \boxed{
 \tau(\Pi_{\ker D_3}-\Pi_{\ker D_4})
 =\frac{\operatorname{rank}\Pi_{\ker D_3}
       -\operatorname{rank}\Pi_{\ker D_4}}{12}
 =\frac{3-4}{12}
 =-\frac1{12}.}
\]
This realization expresses the rational as the transverse defect of steps
three and four in the twelve-phase cycle.

\begin{remark}[Domains of the realizations]
The APP defect \(54\), the eta exponent \(-1/12\), the combinatorial
defect of the hexad--octad lift, and the
modular identities once \(t_3\) has been fixed are exact. The value
of the pseudoinverse, the normalized extractor, and the projector defect are likewise exact.
Each equality preserves its domain and its corresponding map.
\end{remark}
